%!TEX program = lualatex
% Arquivo: relatorio_pratica_6_alcalinidade_agua_mar.tex

\documentclass[a4paper,12pt]{article}

% =========================================================
% Idioma e tipografia
% =========================================================
\usepackage[brazilian]{babel}
\usepackage{fontspec}
\setmainfont{Latin Modern Roman}
\usepackage{microtype}

% =========================================================
% Layout
% =========================================================
\usepackage[top=3cm,bottom=2.5cm,left=3cm,right=2.5cm]{geometry}
\usepackage{setspace}
\onehalfspacing
\usepackage{indentfirst}
\setlength{\parindent}{1.25cm}

% =========================================================
% Matemática, química e unidades
% =========================================================
\usepackage{amsmath,amssymb}
\usepackage[version=4]{mhchem}
\usepackage{siunitx}
\sisetup{
  locale = DE,
  output-decimal-marker = {,},
  per-mode = symbol,
  exponent-product = \cdot
}

% =========================================================
% Tabelas, figuras e listas
% =========================================================
\usepackage{booktabs}
\usepackage{array}
\usepackage{graphicx}
\usepackage{float}
\usepackage{enumitem}
\usepackage{xcolor}
\usepackage{tikz}
\usetikzlibrary{arrows.meta,positioning,calc}
\usepackage[most]{tcolorbox}

\setlist[itemize]{noitemsep,topsep=3pt}
\setlist[enumerate]{noitemsep,topsep=3pt}

% =========================================================
% Links
% =========================================================
\usepackage{hyperref}
\hypersetup{
  colorlinks=true,
  linkcolor=black,
  urlcolor=blue,
  citecolor=black,
  pdfauthor={Israel Carneiro; José Mauro Xavier Elsas; Miguel Gayani},
  pdftitle={Prática 6 -- Determinação da alcalinidade em água do mar}
}

% =========================================================
% Dados
% =========================================================
\newcommand{\ConcHClEstoque}{0.0869}
\newcommand{\VolHClEstoque}{25.0}
\newcommand{\VolHClFinal}{250.0}
\newcommand{\ConcHClTrabalho}{0.00869}
\newcommand{\VolAmostra}{50.00}

\begin{document}

% =========================================================
% Capa
% =========================================================
\begin{titlepage}
\centering

{\large Universidade do Estado do Rio de Janeiro -- UERJ\par}
\vspace{0.3cm}
{\large Química VIII / Oceanografia\par}

\vspace{1.2cm}

{\LARGE\bfseries Prática 6\par}
\vspace{0.4cm}
{\Large\bfseries DETERMINAÇÃO DA ALCALINIDADE\\
EM ÁGUA DO MAR\par}

\vspace{0.8cm}

\IfFileExists{uerj.jpg}{
  \includegraphics[width=0.28\linewidth]{uerj.jpg}
}{
  \begin{tcolorbox}[
    width=0.30\textwidth,
    colback=white,
    colframe=black!35,
    boxrule=0.5pt,
    arc=0mm
  ]
  \centering
  \vspace{1.0cm}
  \textit{Logotipo da UERJ}
  \vspace{1.0cm}
  \end{tcolorbox}
}

\vfill

\begin{flushright}
\textbf{Alunos:}\\
Israel Carneiro\\
José Mauro Xavier Elsas\\
Miguel Gayani\\[0.5cm]
\textbf{Professor:} Ítalo Vieira
\end{flushright}

\vspace{1cm}
Rio de Janeiro -- 2026

\end{titlepage}

% =========================================================
\section{Observação sobre a numeração da prática}

Na apostila da disciplina, a atividade \textit{Determinação da alcalinidade
em água do mar} aparece como \textbf{Prática 5}. Entretanto, por razões de
organização do laboratório, houve uma inversão na sequência das atividades,
e este experimento foi efetivamente realizado como \textbf{Prática 6}.
Neste relatório será adotada a numeração correspondente à ordem em que a
atividade foi executada.

% =========================================================
\section{Introdução}

A alcalinidade de uma água representa sua capacidade de neutralizar íons
\ce{H3O+}, ou, de forma equivalente, sua capacidade de aceitar prótons.
Nas águas naturais, essa propriedade está associada principalmente às
espécies do sistema carbonato, sobretudo bicarbonato
(\ce{HCO3-}) e carbonato (\ce{CO3^2-}), além de possíveis contribuições de
hidróxido (\ce{OH-}) e, em menores proporções, de boratos, fosfatos,
silicatos e espécies orgânicas de caráter básico.

O dióxido de carbono dissolvido participa de equilíbrios ácido--base que
originam as espécies carbonatadas:

\begin{align}
\ce{CO2(g) + 2H2O(l)}
&\ce{<=> H3O+(aq) + HCO3-(aq)},\\
\ce{HCO3-(aq) + H2O(l)}
&\ce{<=> H3O+(aq) + CO3^2-(aq)}.
\end{align}

A determinação da alcalinidade foi realizada por titulação de neutralização
com solução ácida. Foram empregados dois indicadores para identificar
diferentes pontos de neutralização.

Na primeira determinação, utilizou-se fenolftaleína. A viragem desse
indicador permite avaliar a denominada \textit{alcalinidade à
fenolftaleína}, associada à neutralização de \ce{OH-} e à conversão de
\ce{CO3^2-} em \ce{HCO3-}:

\begin{align}
\ce{OH-(aq) + H3O+(aq)}
&\ce{-> 2H2O(l)},\\
\ce{CO3^2-(aq) + H3O+(aq)}
&\ce{-> HCO3-(aq) + H2O(l)}.
\end{align}

Na segunda determinação, empregou-se alaranjado de metila para obtenção da
\textit{alcalinidade total}. Nesse ponto, também ocorre a neutralização do
bicarbonato:

\begin{equation}
\ce{HCO3-(aq) + H3O+(aq) -> CO2(g) + 2H2O(l)}.
\end{equation}

% =========================================================
\section{Objetivos}

\begin{itemize}
  \item preparar, por diluição, a solução de HCl utilizada como titulante;
  \item determinar a alcalinidade à fenolftaleína de uma amostra de água do mar;
  \item determinar a alcalinidade total da amostra com alaranjado de metila;
  \item expressar os resultados em \si{\mole\per\litre} de \ce{H3O+} e em
        \si{\milli\gram\per\litre} como \ce{CaCO3};
  \item estimar as contribuições relativas de hidróxido, carbonato e
        bicarbonato para a alcalinidade da amostra.
\end{itemize}

% =========================================================
\section{Material utilizado}

\begin{itemize}
  \item bureta de \qty{50.00}{\milli\litre};
  \item Erlenmeyers de \qty{250}{\milli\litre};
  \item balão volumétrico de \qty{250.00}{\milli\litre};
  \item pipeta para transferência de \qty{25.0}{\milli\litre};
  \item material volumétrico para transferência de
        \qty{50.00}{\milli\litre} de amostra;
  \item suporte universal e garra para bureta;
  \item frasco lavador com água destilada.
\end{itemize}

% =========================================================
\section{Reagentes e soluções}

\begin{itemize}
  \item solução de HCl previamente padronizada, de concentração
        \qty{\ConcHClEstoque}{\mole\per\litre};
  \item solução diluída de HCl usada como titulante;
  \item fenolftaleína;
  \item alaranjado de metila;
  \item amostra de água do mar;
  \item água destilada.
\end{itemize}

% =========================================================
\section{Procedimento experimental}

\subsection{Preparo da solução diluída de HCl}

O quadro da aula indicava o preparo de uma solução aproximadamente
\qty{0.01}{\mole\per\litre} por diluição de \qty{25.0}{\milli\litre} da
solução de HCl para um volume final de \qty{250.0}{\milli\litre}.
Como a solução estoque havia sido previamente padronizada e sua concentração
experimental era \qty{\ConcHClEstoque}{\mole\per\litre}, utilizou-se esse
valor exato no cálculo da concentração efetiva da solução diluída.

A relação de diluição é:

\begin{equation}
C_1V_1=C_2V_2.
\end{equation}

Assim:

\begin{align}
C_2
&=
\frac{C_1V_1}{V_2}\\
&=
\frac{
\left(\qty{\ConcHClEstoque}{\mole\per\litre}\right)
\left(\qty{\VolHClEstoque}{\milli\litre}\right)}
{\qty{\VolHClFinal}{\milli\litre}}\\
&=
\boxed{\qty{\ConcHClTrabalho}{\mole\per\litre}}.
\end{align}

Portanto, embora a solução seja referida no roteiro como HCl
\qty{0.010}{\mole\per\litre}, a concentração efetivamente utilizada nos
cálculos deste relatório foi \qty{\ConcHClTrabalho}{\mole\per\litre}.

\subsection{Alcalinidade à fenolftaleína}

Foram transferidos \qty{\VolAmostra}{\milli\litre} da amostra de água do
mar para um Erlenmeyer de \qty{250}{\milli\litre}. Adicionou-se
fenolftaleína e observou-se a coloração da amostra.

Durante a execução ocorreu uma dificuldade experimental importante: foram
necessárias \textbf{oito gotas de fenolftaleína} para que a solução
apresentasse uma coloração apenas levemente rosada. Em seguida, a amostra
foi titulada com a solução de HCl até o desaparecimento da coloração rosa.
O procedimento foi realizado em triplicata.

\subsection{Alcalinidade total com alaranjado de metila}

Em outra alíquota de \qty{\VolAmostra}{\milli\litre} da mesma amostra,
adicionou-se alaranjado de metila conforme a orientação experimental do
quadro. A solução foi titulada com HCl até a mudança da coloração amarela
para vermelho-rosa. A determinação também foi realizada em triplicata.

% =========================================================
\section{Resultados e cálculos}

\subsection{Volumes de HCl consumidos}

Os volumes experimentais obtidos foram:

\begin{table}[H]
\centering
\caption{Volumes de HCl consumidos nas determinações de alcalinidade.}
\label{tab:volumes}
\begin{tabular}{
c
S[table-format=2.2]
S[table-format=2.1]
}
\toprule
\textbf{Determinação}
& {\textbf{Fenolftaleína} (\si{\milli\litre})}
& {\textbf{Alaranjado de metila} (\si{\milli\litre})}\\
\midrule
1 & 0,60 & 12,7\\
2 & 0,60 & 12,4\\
3 & 0,50 & 12,6\\
\midrule
\textbf{Média} & 0,57 & 12,57\\
\bottomrule
\end{tabular}
\end{table}

Para a fenolftaleína:

\begin{align}
\overline{V}_{P}
&=
\frac{0.60+0.60+0.50}{3}\\
&=
\qty{0.5667}{\milli\litre}
\approx
\boxed{\qty{0.57}{\milli\litre}}.
\end{align}

Para a alcalinidade total:

\begin{align}
\overline{V}_{T}
&=
\frac{12.7+12.4+12.6}{3}\\
&=
\boxed{\qty{12.57}{\milli\litre}}.
\end{align}

\subsection{Alcalinidade em mol L\texorpdfstring{$^{-1}$}{-1} de \ce{H3O+}}

A alcalinidade correspondente ao volume de ácido consumido pode ser
expressa pela razão entre a quantidade de matéria de \ce{H3O+} fornecida
pelo titulante e o volume da amostra:

\begin{equation}
A=
\frac{C_{\ce{HCl}}V_{\ce{HCl}}}{V_{\text{amostra}}}.
\label{eq:alcalinidade_molar}
\end{equation}

\subsubsection*{Alcalinidade à fenolftaleína}

\begin{align}
P
&=
\frac{
\left(\qty{0.00869}{\mole\per\litre}\right)
\left(\qty{0.5667e-3}{\litre}\right)}
{\qty{0.05000}{\litre}}\\
&=
\boxed{\qty{9.85e-5}{\mole\per\litre}}.
\end{align}

\subsubsection*{Alcalinidade total}

\begin{align}
T
&=
\frac{
\left(\qty{0.00869}{\mole\per\litre}\right)
\left(\qty{12.5667e-3}{\litre}\right)}
{\qty{0.05000}{\litre}}\\
&=
\boxed{\qty{2.184e-3}{\mole\per\litre}}.
\end{align}

% =========================================================
\subsection{Alcalinidade expressa como \ce{CaCO3}}

Para expressar a alcalinidade como \ce{CaCO3}, empregou-se o fator
\qty{50000}{\milli\gram\per\equiv}, correspondente à massa equivalente do
\ce{CaCO3} expressa em miligramas. Assim:

\begin{equation}
A\left(\si{\milli\gram\per\litre}\ \text{como \ce{CaCO3}}\right)
=
A\left(\si{\equiv\per\litre}\right)
\times
50000.
\end{equation}

Como o HCl é monoprótico, a concentração em
\si{\mole\per\litre} de \ce{H3O+} é numericamente igual à concentração
em \si{\equiv\per\litre}.

Para a alcalinidade à fenolftaleína:

\begin{align}
P
&=
\left(9.85\times10^{-5}\right)
\left(50000\right)\\
&=
\boxed{\qty{4.92}{\milli\gram\per\litre}\ \text{como \ce{CaCO3}}}.
\end{align}

Para a alcalinidade total:

\begin{align}
T
&=
\left(2.184\times10^{-3}\right)
\left(50000\right)\\
&=
\boxed{\qty{109.2}{\milli\gram\per\litre}\ \text{como \ce{CaCO3}}}.
\end{align}

\subsection{Resultados individuais}

\begin{table}[H]
\centering
\caption{Alcalinidade calculada para cada determinação.}
\label{tab:resultados_individuais}
\begin{tabular}{
c
S[table-format=3.3]
S[table-format=3.3]
}
\toprule
\textbf{Determinação}
& {\textbf{$P$} (\si{\milli\gram\per\litre} como \ce{CaCO3})}
& {\textbf{$T$} (\si{\milli\gram\per\litre} como \ce{CaCO3})}\\
\midrule
1 & 5,214 & 110,363\\
2 & 5,214 & 107,756\\
3 & 4,345 & 109,494\\
\midrule
\textbf{Média} & 4,924 & 109,204\\
\bottomrule
\end{tabular}
\end{table}

% =========================================================
\subsection{Componentes da alcalinidade}

A apostila relaciona a alcalinidade à fenolftaleína ($P$) e a
alcalinidade total ($T$) com as contribuições de \ce{OH-},
\ce{CO3^2-} e \ce{HCO3-}.

Neste experimento:

\begin{equation}
P=\qty{4.92}{\milli\gram\per\litre}
\qquad\text{e}\qquad
T=\qty{109.2}{\milli\gram\per\litre},
\end{equation}

ambos expressos como \ce{CaCO3}. Como:

\begin{equation}
P<\frac{T}{2},
\end{equation}

aplica-se a condição:

\begin{align}
\text{hidróxido} &= 0,\\
\text{carbonato} &= 2P,\\
\text{bicarbonato} &= T-2P.
\end{align}

Portanto:

\begin{align}
\text{Contribuição de \ce{OH-}}
&=
\boxed{0},\\[0.2cm]
\text{Contribuição de \ce{CO3^2-}}
&=
2(4.92)\\
&=
\boxed{\qty{9.84}{\milli\gram\per\litre}\ \text{como \ce{CaCO3}}},\\[0.2cm]
\text{Contribuição de \ce{HCO3-}}
&=
109.2-9.84\\
&=
\boxed{\qty{99.4}{\milli\gram\per\litre}\ \text{como \ce{CaCO3}}}.
\end{align}

Esses valores representam \textbf{contribuições para a alcalinidade
expressas como \ce{CaCO3}}, e não as concentrações mássicas reais dos íons
\ce{CO3^2-} e \ce{HCO3-}.

% =========================================================
\section{Discussão}

A alcalinidade total determinada para a amostra foi de aproximadamente
\qty{109.2}{\milli\gram\per\litre} como \ce{CaCO3}. De acordo com a
classificação apresentada na apostila da disciplina, esse valor situa a
amostra na faixa de alcalinidade média.

A alcalinidade à fenolftaleína foi muito inferior à alcalinidade total:

\begin{equation}
P=\qty{4.92}{\milli\gram\per\litre}
\ll
\frac{T}{2}
=
\qty{54.6}{\milli\gram\per\litre}.
\end{equation}

Essa relação conduz, segundo o critério apresentado no roteiro, à ausência
de contribuição calculada de \ce{OH-}, à presença de uma pequena parcela
atribuída a \ce{CO3^2-} e à predominância da contribuição de
\ce{HCO3-}.

Um aspecto importante observado durante a determinação com fenolftaleína
foi a dificuldade de visualização da coloração inicial. Foram necessárias
oito gotas do indicador para que a amostra apresentasse apenas uma
coloração levemente rosada. Esse comportamento indica que a amostra se
encontrava muito próxima da faixa de transição do indicador e é coerente
com o pequeno volume de HCl consumido nessa etapa.

Os volumes medidos na titulação à fenolftaleína foram apenas
\qty{0.60}{\milli\litre}, \qty{0.60}{\milli\litre} e
\qty{0.50}{\milli\litre}. Para volumes tão pequenos, uma diferença de uma
gota de titulante ou uma pequena diferença na percepção visual da viragem
representa uma fração relativamente elevada do resultado. Por isso, a
alcalinidade à fenolftaleína apresenta maior incerteza relativa que a
alcalinidade total.

Essa limitação, entretanto, não altera a interpretação geral dos
componentes da alcalinidade, pois mesmo uma variação relativamente grande
em $P$ não seria suficiente para aproximá-lo de $T/2$. Assim, a condição
$P<T/2$ permanece bem estabelecida pelos dados experimentais.

Na determinação com alaranjado de metila, os volumes de
\qty{12.7}{\milli\litre}, \qty{12.4}{\milli\litre} e
\qty{12.6}{\milli\litre} mostraram maior concordância entre as replicatas,
tornando a determinação da alcalinidade total mais robusta.

Também é importante observar que a solução de HCl utilizada não possuía
exatamente a concentração nominal de \qty{0.010}{\mole\per\litre}.
Como foi preparada a partir da solução previamente padronizada de
\qty{0.0869}{\mole\per\litre}, utilizando-se
\qty{25.0}{\milli\litre} e diluindo-se para
\qty{250.0}{\milli\litre}, sua concentração efetiva foi
\qty{0.00869}{\mole\per\litre}. O uso desse valor é necessário para que os
resultados reflitam as condições efetivamente empregadas no laboratório.

% =========================================================
\section{Conclusão}

A determinação da alcalinidade da amostra de água do mar foi realizada por
titulação com solução de HCl de concentração efetiva
\qty{0.00869}{\mole\per\litre}.

A alcalinidade à fenolftaleína obtida foi:

\begin{equation*}
\boxed{
P=
\qty{9.85e-5}{\mole\per\litre}\ \text{de \ce{H3O+}}
=
\qty{4.92}{\milli\gram\per\litre}\ \text{como \ce{CaCO3}}
}.
\end{equation*}

A alcalinidade total foi:

\begin{equation*}
\boxed{
T=
\qty{2.184e-3}{\mole\per\litre}\ \text{de \ce{H3O+}}
=
\qty{109.2}{\milli\gram\per\litre}\ \text{como \ce{CaCO3}}
}.
\end{equation*}

Como $P<T/2$, a interpretação dos componentes indicou ausência de
contribuição calculada de hidróxido, uma contribuição de aproximadamente
\qty{9.84}{\milli\gram\per\litre} atribuída ao carbonato e uma contribuição
predominante de aproximadamente
\qty{99.4}{\milli\gram\per\litre} atribuída ao bicarbonato, todas expressas
como \ce{CaCO3}.

A determinação à fenolftaleína deve ser interpretada com maior cautela,
pois foram necessárias oito gotas de indicador para produzir uma coloração
apenas levemente rosada e os volumes de titulante consumidos foram muito
pequenos. Apesar dessa limitação experimental, a relação $P<T/2$ e a
predominância da contribuição bicarbonato permaneceram claramente
definidas.

% =========================================================
\section*{Referência}

LAGO, Dalva Cristina Baptista do; GOIS, Jefferson Santos de; LANGONE,
Marta Antunes Pereira; SENNA, Lilian Ferreira de.
\textit{Química VIII -- Oceanografia: aulas práticas}.
Rio de Janeiro: Universidade do Estado do Rio de Janeiro,
Instituto de Química, Departamento de Química Analítica, 2017.

\end{document}
